% =====================================================================
%  phuluc_nguon.tex  —  PHỤ LỤC: BẢNG ĐỐI CHIẾU LUẬN ĐIỂM - NGUỒN TRÍCH DẪN
% =====================================================================

\chapter*{PHỤ LỤC A. ĐỐI CHIẾU LUẬN ĐIỂM VÀ NGUỒN TRÍCH DẪN}
\addcontentsline{toc}{chapter}{PHỤ LỤC A. ĐỐI CHIẾU LUẬN ĐIỂM VÀ NGUỒN TRÍCH DẪN}
\label{app:nguon}

Phụ lục này liệt kê từng luận điểm lý thuyết và phương pháp luận được sử dụng trong dự án kèm nguồn gốc học thuật tương ứng, nhằm bảo đảm mọi khẳng định đều truy nguyên được về tài liệu công bố chính thống.

\begin{longtable}{|p{0.08\textwidth}|p{0.48\textwidth}|p{0.34\textwidth}|}
\hline
\textbf{Mục} & \textbf{Luận điểm} & \textbf{Nguồn} \\
\hline
\endfirsthead
\hline
\textbf{Mục} & \textbf{Luận điểm} & \textbf{Nguồn} \\
\hline
\endhead

\hline
\multicolumn{3}{|l|}{\textit{A.1. Tổng quan về Hệ thống gợi ý \& Lọc cộng tác}} \\
\hline
1.1.1 & Vai trò cá nhân hóa gợi ý trong TMĐT/chuỗi cung ứng &
\textcite{schafer2001commerce}; \textcite{linden2003amazon} \\
\hline
2.1.2 & Nền tảng Lọc cộng tác dựa trên bộ nhớ (Memory-based CF, GroupLens, Item-based CF) &
\textcite{resnick1994grouplens}; \textcite{sarwar2001item}; \textcite{su2009survey} \\
\hline
1.1.2 & Hạn chế của tích vô hướng trong Nhân tử hóa ma trận (Matrix Factorization) &
\textcite{koren2009matrix} \\
\hline
1.2.2 & Bộ dữ liệu DataCo Smart Supply Chain Dataset &
\textcite{dataco2019} \\
\hline

\hline
\multicolumn{3}{|l|}{\textit{A.2. Phản hồi ẩn \& Kiến trúc Neural Collaborative Filtering}} \\
\hline
2.2.4 & Xếp hạng cá nhân hóa Bayes (BPR) từ phản hồi ẩn &
\textcite{rendle2009bpr} \\
\hline
2.3 & Khung làm việc NCF và kiến trúc mô hình lai NeuMF (GMF + MLP) &
\textcite{he2017ncf} \\
\hline
2.3.4 & Chiến lược kết hợp Early Fusion và Late Fusion &
\textcite{he2017ncf} \\
\hline
3.6 & Thuật toán tối ưu hóa thích ứng Adam (Pre-training GMF/MLP) &
\textcite{kingma2014adam} \\
\hline
2.4.1 & Mô hình Wide \& Deep và DeepFM trong gợi ý (Feature-rich Deep CTR) &
\textcite{cheng2016wide}; \textcite{guo2017deepfm} \\
\hline

\hline
\multicolumn{3}{|l|}{\textit{A.3. Xu hướng nghiên cứu hiện đại (2016--2025)}} \\
\hline
2.4.2 & Kiến trúc hai tháp (Two-Tower) cho truy hồi quy mô lớn &
\textcite{covington2016youtube} \\
\hline
2.4.3 & Gợi ý theo chuỗi hành vi (Sequential/Session-based) &
\textcite{hidasi2016gru4rec}; \textcite{kang2018self}; \textcite{sun2019bert4rec} \\
\hline
2.4.4 & Học biểu diễn trên đồ thị (Graph-based Collaborative Filtering) &
\textcite{wang2019ngcf}; \textcite{he2020lightgcn} \\
\hline
2.4.5 & Biến phân tự mã hóa, học tự giám sát, LLM và Diffusion cho gợi ý &
\textcite{liang2018multvae}; \textcite{wu2021sgl}; \textcite{rajput2023tiger}; \textcite{wu2023survey_llm4rec}; \textcite{wei2025diffusion} \\
\hline
2.4 & Khảo sát cập nhật (2024) về mạng nơ-ron sâu trong Collaborative Filtering &
\textcite{li2024dnncf}; \textcite{zhang2019deeplearningsurvey} \\
\hline

\hline
\multicolumn{3}{|l|}{\textit{A.4. Đánh giá Top-K, Cài đặt \& Hướng phát triển}} \\
\hline
2.5 & Giao thức đánh giá hiệu năng xếp hạng Top-N Recommendation &
\textcite{cremonesi2010performance}; \textcite{he2017ncf} \\
\hline
4.1 & Vấn đề đánh giá bằng mẫu âm lấy mẫu con (Sampled Metrics) &
\textcite{krichene2020sampled} \\
\hline
1.3 & Nền tảng và thư viện tính toán học sâu PyTorch &
\textcite{paszke2019pytorch} \\
\hline
2.5 & Đánh giá xếp hạng dựa trên độ đo NDCG &
\textcite{jarvelin2002} \\
\hline
2.5.1 & Chỉ số ngoài độ chính xác: Novelty (Self-Information) &
\textcite{vargas2011novelty} \\
\hline
\end{longtable}

\vspace{0.5\baselineskip}
\noindent\textbf{Ghi chú.} Các thông số cấu hình và thiết kế hệ thống trình bày trong Chương~\ref{chap:thietke_trienkhai} và kết quả thực nghiệm trong Chương~\ref{chap:thucnghiem} do nhóm tác giả tự triển khai và tính toán từ mã nguồn thực nghiệm của dự án.